\documentclass[11pt]{article}

\usepackage[margin=1in]{geometry}
\usepackage{amsmath,amssymb,amsthm}
\usepackage{enumitem}
\usepackage{mathtools}

\setlength{\parindent}{0pt}
\setlength{\parskip}{0.45em}

\title{Homework 1 -- Theory}
\author{}
\date{}

\begin{document}
\maketitle

% ============================================================
\section*{1}

Let \(D\) be any distribution on \(\{-1,1\}^n\), and let
\[
f,g:\{-1,1\}^n\to\{-1,1\}.
\]

\subsection*{(a)}

Since \(f(x),g(x)\in\{-1,1\}\),
\[
f(x)g(x)
=
\begin{cases}
+1, & f(x)=g(x),\\
-1, & f(x)\neq g(x).
\end{cases}
\]

Let
\[
p=\Pr_{x\sim D}[f(x)\neq g(x)].
\]
Then
\[
\Pr_{x\sim D}[f(x)=g(x)]=1-p.
\]
Therefore
\begin{align*}
\mathbb{E}_{x\sim D}[f(x)g(x)]
&=(+1)\Pr[f(x)=g(x)]+(-1)\Pr[f(x)\neq g(x)]\\
&=(1-p)-p\\
&=1-2p.
\end{align*}
Solving for \(p\),
\[
2p=1-\mathbb{E}_{x\sim D}[f(x)g(x)],
\]
so
\[
\boxed{
\Pr_{x\sim D}[f(x)\neq g(x)]
=
\frac{1-\mathbb{E}_{x\sim D}[f(x)g(x)]}{2}.
}
\]

\subsection*{(b)}

Yes. The argument only uses that \(f(x)\) and \(g(x)\) take values in
\(\{-1,1\}\). It does not depend on the domain being \(\{-1,1\}^n\).
Also, \(|f(x)g(x)|=1\), so the expectation is defined for any distribution
\(D\). Thus the same identity holds for any domain, including
\(\mathbb{R}^n\).

% ============================================================
\section*{2}

Let \(f\) be a decision tree with \(t\) leaves over
\[
x=(x_1,\ldots,x_n)\in\{-1,1\}^n.
\]
Interpret \(1\) as TRUE and \(-1\) as FALSE.

For a leaf \(\ell\), suppose the path to that leaf requires
\[
x_{i_{\ell,1}}=b_{\ell,1},\ldots,
x_{i_{\ell,d_\ell}}=b_{\ell,d_\ell},
\qquad
b_{\ell,j}\in\{-1,1\},
\]
and let the label at the leaf be
\[
y_\ell\in\{-1,1\}.
\]

For one condition \(x_i=b\), consider
\[
\frac{1+bx_i}{2}.
\]
Since \(b^2=1\),
\[
x_i=b
\quad\Longrightarrow\quad
\frac{1+bx_i}{2}
=
\frac{1+b^2}{2}
=
1,
\]
while
\[
x_i=-b
\quad\Longrightarrow\quad
\frac{1+bx_i}{2}
=
\frac{1-b^2}{2}
=
0.
\]
Thus
\[
\frac{1+bx_i}{2}
=
\begin{cases}
1, & x_i=b,\\
0, & x_i\neq b.
\end{cases}
\]

Multiplying these factors along the path gives the indicator polynomial
\[
I_\ell(x)
=
\prod_{j=1}^{d_\ell}
\frac{1+b_{\ell,j}x_{i_{\ell,j}}}{2}.
\]
Each factor is \(0\) or \(1\), so \(I_\ell(x)=1\) exactly when \(x\)
follows the path to leaf \(\ell\), and \(I_\ell(x)=0\) otherwise.

Define
\[
\boxed{
p(x)
=
\sum_{\ell=1}^{t}y_\ell I_\ell(x)
=
\sum_{\ell=1}^{t}
y_\ell
\prod_{j=1}^{d_\ell}
\frac{1+b_{\ell,j}x_{i_{\ell,j}}}{2}.
}
\]

For any \(x\in\{-1,1\}^n\), exactly one leaf is reached, since at each
internal node the queried variable has a single value and one branch is
taken. If that leaf is \(\ell(x)\), then
\[
I_{\ell(x)}(x)=1.
\]
Any other leaf's path starts at the same root but ends elsewhere, so it
diverges at some node and requires a value that \(x\) does not have. That
factor is \(0\), so its indicator is \(0\). Hence
\[
p(x)
=
y_{\ell(x)}\cdot 1
+
\sum_{\ell\neq\ell(x)}y_\ell\cdot 0
=
y_{\ell(x)}
=
f(x).
\]
Therefore
\[
\boxed{
p(x)=f(x)
\qquad
\text{for every }x\in\{-1,1\}^n.
}
\]

% ============================================================
\section*{3}

The Gini function is
\[
C(a)=2a(1-a),
\]
where \(a\) is the fraction of negative examples at a node. For a split on
a variable \(X_i\),
\[
\mathrm{new}(X_i)
=
\Pr_S[X_i=0]\,
C\!\left(\Pr_S[\mathrm{neg}\mid X_i=0]\right)
+
\Pr_S[X_i=1]\,
C\!\left(\Pr_S[\mathrm{neg}\mid X_i=1]\right),
\]
where \(\mathrm{old}=C(\Pr_S[\mathrm{neg}])\) is the impurity at the node
before the split, and
\[
\mathrm{gain}(X_i)
=
\mathrm{old}-\mathrm{new}(X_i).
\]
The tree is required to be a full binary tree of depth two, so every internal
node is split even if both of its children carry the same majority label.

\paragraph{Root counts.}

The total number of positive examples is
\[
10+25+35+35+5+30+10+15=165,
\]
and the total number of negative examples is
\[
20+5+15+5+15+10+10+5=85.
\]
Thus
\[
|S|=250,
\]
and the impurity before the first split is
\[
\mathrm{old}
=
C\!\left(\frac{85}{250}\right)
=
2\left(\frac{85}{250}\right)\left(\frac{165}{250}\right)
=
0.4488.
\]

\paragraph{Root split on \(X\).}

For \(X=0\),
\[
105\text{ pos},\qquad45\text{ neg},\qquad150\text{ total},
\]
so
\[
C\!\left(\frac{45}{150}\right)
=
2\left(\frac{45}{150}\right)\left(\frac{105}{150}\right)
=
0.42.
\]
For \(X=1\),
\[
60\text{ pos},\qquad40\text{ neg},\qquad100\text{ total},
\]
so
\[
C\!\left(\frac{40}{100}\right)
=
2(0.4)(0.6)
=
0.48.
\]
Therefore
\[
\mathrm{new}(X)
=
\frac{150}{250}(0.42)
+
\frac{100}{250}(0.48)
=
0.444,
\]
and
\[
\mathrm{gain}(X)=0.4488-0.444=0.0048.
\]

\paragraph{Root split on \(Y\).}

For \(Y=0\),
\[
70\text{ pos},\qquad50\text{ neg},\qquad120\text{ total},
\]
so
\[
C\!\left(\frac{50}{120}\right)
=
2\left(\frac{5}{12}\right)\left(\frac{7}{12}\right)
=
\frac{35}{72}
\approx0.4861.
\]
For \(Y=1\),
\[
95\text{ pos},\qquad35\text{ neg},\qquad130\text{ total},
\]
so
\[
C\!\left(\frac{35}{130}\right)
=
2\left(\frac{7}{26}\right)\left(\frac{19}{26}\right)
=
\frac{133}{338}
\approx0.3935.
\]
Hence
\[
\mathrm{new}(Y)
=
\frac{120}{250}\cdot\frac{35}{72}
+
\frac{130}{250}\cdot\frac{133}{338}
=
\frac{427}{975}
\approx0.4379,
\]
and
\[
\mathrm{gain}(Y)
=
0.4488-0.4379
\approx0.0109.
\]

\paragraph{Root split on \(Z\).}

For \(Z=0\),
\[
60\text{ pos},\qquad60\text{ neg},\qquad120\text{ total},
\]
so
\[
C\!\left(\frac{60}{120}\right)=0.5.
\]
For \(Z=1\),
\[
105\text{ pos},\qquad25\text{ neg},\qquad130\text{ total},
\]
so
\[
C\!\left(\frac{25}{130}\right)
=
2\left(\frac{5}{26}\right)\left(\frac{21}{26}\right)
=
\frac{105}{338}
\approx0.3107.
\]
Thus
\[
\mathrm{new}(Z)
=
\frac{120}{250}(0.5)
+
\frac{130}{250}\cdot\frac{105}{338}
=
\frac{261}{650}
\approx0.4015,
\]
and
\[
\mathrm{gain}(Z)
=
0.4488-0.4015
\approx0.0473.
\]

Comparing the gains,
\[
0.0473>0.0109>0.0048,
\]
so the root split is
\[
\boxed{Z}.
\]

\paragraph{Node \(Z=0\).}

This node contains
\[
60\text{ pos},\qquad60\text{ neg},
\]
so
\[
\mathrm{old}
=
C\!\left(\frac{60}{120}\right)
=
0.5.
\]

For a split on \(X\),
\[
X=0:\quad45\text{ pos},\ 35\text{ neg},\ 80\text{ total},
\]
with
\[
C\!\left(\frac{35}{80}\right)
=
\frac{63}{128}
\approx0.4922,
\]
and
\[
X=1:\quad15\text{ pos},\ 25\text{ neg},\ 40\text{ total},
\]
with
\[
C\!\left(\frac{25}{40}\right)
=
\frac{15}{32}
=
0.46875.
\]
Therefore
\[
\mathrm{new}(X)
=
\frac{80}{120}\cdot\frac{63}{128}
+
\frac{40}{120}\cdot\frac{15}{32}
=
\frac{31}{64}
=
0.484375,
\]
so
\[
\mathrm{gain}(X)
=
0.5-0.484375
=
0.015625.
\]

For a split on \(Y\),
\[
Y=0:\quad15\text{ pos},\ 35\text{ neg},\ 50\text{ total},
\]
with
\[
C\!\left(\frac{35}{50}\right)=0.42,
\]
and
\[
Y=1:\quad45\text{ pos},\ 25\text{ neg},\ 70\text{ total},
\]
with
\[
C\!\left(\frac{25}{70}\right)
=
\frac{45}{98}
\approx0.4592.
\]
Therefore
\[
\mathrm{new}(Y)
=
\frac{50}{120}(0.42)
+
\frac{70}{120}\cdot\frac{45}{98}
=
\frac{31}{70}
\approx0.4429,
\]
so
\[
\mathrm{gain}(Y)
=
0.5-\frac{31}{70}
=
\frac{2}{35}
\approx0.0571.
\]

Since
\[
0.0571>0.015625,
\]
this node splits on \(Y\). The majority labels are
\[
Z=0,\ Y=0:
\quad15\text{ pos},\ 35\text{ neg}
\Rightarrow -1,
\]
and
\[
Z=0,\ Y=1:
\quad45\text{ pos},\ 25\text{ neg}
\Rightarrow +1.
\]

\paragraph{Node \(Z=1\).}

This node contains
\[
105\text{ pos},\qquad25\text{ neg},
\]
so
\[
\mathrm{old}
=
C\!\left(\frac{25}{130}\right)
=
\frac{105}{338}
\approx0.3107.
\]

For a split on \(X\),
\[
X=0:\quad60\text{ pos},\ 10\text{ neg},\ 70\text{ total},
\]
with
\[
C\!\left(\frac{10}{70}\right)
=
\frac{12}{49}
\approx0.2449,
\]
and
\[
X=1:\quad45\text{ pos},\ 15\text{ neg},\ 60\text{ total},
\]
with
\[
C\!\left(\frac{15}{60}\right)
=
\frac{3}{8}
=
0.375.
\]
Therefore
\[
\mathrm{new}(X)
=
\frac{70}{130}\cdot\frac{12}{49}
+
\frac{60}{130}\cdot\frac{3}{8}
=
\frac{111}{364}
\approx0.3049,
\]
so
\[
\mathrm{gain}(X)
=
0.3107-0.3049
\approx0.0057.
\]

For a split on \(Y\),
\[
Y=0:\quad55\text{ pos},\ 15\text{ neg},\ 70\text{ total},
\]
with
\[
C\!\left(\frac{15}{70}\right)
=
\frac{33}{98}
\approx0.3367,
\]
and
\[
Y=1:\quad50\text{ pos},\ 10\text{ neg},\ 60\text{ total},
\]
with
\[
C\!\left(\frac{10}{60}\right)
=
\frac{5}{18}
\approx0.2778.
\]
Therefore
\[
\mathrm{new}(Y)
=
\frac{70}{130}\cdot\frac{33}{98}
+
\frac{60}{130}\cdot\frac{5}{18}
=
\frac{13}{42}
\approx0.3095,
\]
so
\[
\mathrm{gain}(Y)
=
0.3107-0.3095
\approx0.0011.
\]

Since
\[
0.0057>0.0011,
\]
this node splits on \(X\). The majority labels are
\[
Z=1,\ X=0:
\quad60\text{ pos},\ 10\text{ neg}
\Rightarrow +1,
\]
and
\[
Z=1,\ X=1:
\quad45\text{ pos},\ 15\text{ neg}
\Rightarrow +1.
\]

The full depth-two tree is
\[
\boxed{
\begin{array}{c}
\text{root: }Z\\[2mm]
Z=0\Rightarrow Y:
\quad Y=0\mapsto -1,\quad Y=1\mapsto +1\\[1mm]
Z=1\Rightarrow X:
\quad X=0\mapsto +1,\quad X=1\mapsto +1.
\end{array}
}
\]

The four leaves classify
\[
35+45+60+45=185
\]
examples correctly and
\[
15+25+10+15=65
\]
incorrectly, with
\[
185+65=250.
\]
Thus
\[
\boxed{
\mathrm{accuracy}
=
\frac{185}{250}
=
0.74
=
74\%.
}
\]

% ============================================================
\section*{4}

Let the target function be \(h_{\theta^*}\). Draw \(m\) labeled examples
\[
S=\{(x^1,y^1),\ldots,(x^m,y^m)\}.
\]

Choose a threshold that is consistent with all of the training examples. If
both labels appear, let
\[
x^-=\max\{x^i:y^i=-1\},
\qquad
x^+=\min\{x^i:y^i=1\},
\]
and choose any
\[
x^-\le\hat\theta<x^+.
\]
If only one label appears, choose any threshold consistent with the sample.
Both \(x^-\) and \(x^+\) are found in one pass through \(S\), so the
algorithm is efficient.

If the chosen threshold has true error greater than \(\epsilon\), then it
must be too far to the left or too far to the right of the true threshold
\(\theta^*\). In either case, the training examples must miss a region with
probability at least \(\epsilon\). Therefore
\[
\Pr[\mathrm{err}(h_{\hat\theta})>\epsilon]
\le
2(1-\epsilon)^m
\le
2e^{-\epsilon m}.
\]

To make this probability at most \(\delta\), require
\[
2e^{-\epsilon m}\le\delta.
\]

Taking logs gives
\[
\boxed{
m\ge\frac{1}{\epsilon}\log\frac{2}{\delta}.
}
\]

Thus
\[
\boxed{
m=
O\!\left(
\frac{1}{\epsilon}\log\frac{1}{\delta}
\right),
}
\]
so the threshold class is PAC learnable with the required number of examples.

% ============================================================
\section*{5}

Assume \(A\) has mistake bound \(t\), and define
\[
\mathrm{err}(h)=\Pr_{x\sim D}[h(x)\neq c(x)].
\]

\subsection*{(a)}

Fix a hypothesis \(h'\) with
\[
\mathrm{err}(h')>\epsilon.
\]

The probability that it correctly classifies one fresh example is less than
\(1-\epsilon\). Therefore, the probability that it correctly classifies a
block of \(k\) fresh examples is at most
\[
(1-\epsilon)^k\le e^{-\epsilon k}.
\]

To make this at most \(\delta'\), choose
\[
e^{-\epsilon k}\le\delta'.
\]

Hence
\[
\boxed{
k\ge\frac{1}{\epsilon}\log\frac{1}{\delta'}.
}
\]

\subsection*{(b)}

Since \(A\) changes its hypothesis only after a mistake and makes at most
\(t\) mistakes, at most \(t+1\) hypotheses need to be considered. Set
\[
\delta'=\frac{\delta}{t+1}.
\]

Using part (a), choose
\[
\boxed{
k=
\left\lceil
\frac{1}{\epsilon}\log\frac{t+1}{\delta}
\right\rceil.
}
\]

Using \(t+1\) blocks gives
\[
\boxed{
m=(t+1)
\left\lceil
\frac{1}{\epsilon}\log\frac{t+1}{\delta}
\right\rceil.
}
\]

\subsection*{(c)}

Algorithm \(B\) uses \(A\) as follows:
\begin{enumerate}[leftmargin=*,itemsep=2pt]
    \item Split the \(m\) examples into \(t+1\) blocks of size \(k\).
    \item Test the current hypothesis of \(A\) on the next block.
    \item If it classifies the whole block correctly, output it.
    \item Otherwise, give \(A\) one misclassified example so that \(A\)
    updates, and continue to the next block.
\end{enumerate}

Every failed block causes one mistake. If all \(t+1\) blocks failed, then
\(A\) would make \(t+1\) mistakes, which contradicts the mistake bound
\(t\). Therefore, some hypothesis must pass a block.

For any bad hypothesis with error greater than \(\epsilon\), part (a) shows
that the probability of passing its block is at most
\[
\frac{\delta}{t+1}.
\]

There are at most \(t+1\) hypotheses, so the union bound gives
\[
\Pr[\mathrm{err}(B)>\epsilon]
\le
(t+1)\frac{\delta}{t+1}
=
\delta.
\]

Thus, with probability at least \(1-\delta\), the output has error at most
\(\epsilon\). The total sample size is
\[
\boxed{
m=
O\!\left(
\frac{t+1}{\epsilon}
\log\frac{t+1}{\delta}
\right).
}
\]

\end{document}